\paragraph{Only a total support bound is necessary.}
Fix $C,a\ge1$ and $r\ge2$.  For a prescribed support
$\Omega\subseteq[r]\times[r]$ with
\[
                     |\Omega|\le Cr(\lceil\log_2r\rceil)^a,
\]
the full permanent polynomial has a deterministically
constructible integral constant-free circuit of size and
bit construction time
\[
 \operatorname{poly}(r)\,
       2^{O_{C,a}(r(\log\log r)^2/\sqrt{\log r})},
\]
and formal degree at most $2r$, with the subtraction convention
used above.  Thus no bound on any individual row or column
degree is needed.  The original parameter is still the
matrix order $r$.

Here is the full reduction to the preceding row-scope result.
Let $\lambda=\lceil\log_2r\rceil$, set
$T=\lceil\lambda^{a+2}\rceil$, and let $H$ be the set of
rows containing more than $T$ allowed entries.  If $h=|H|$,
then counting allowed entries gives
\[
                          h\le Cr/\lambda^2.
\]
For sufficiently large $r$, this is at most $r/2$.
Expand the definition of the permanent by first specifying
the columns matched to the rows of $H$.  It gives the exact
integral polynomial identity
\[
 \operatorname{per}(A)=
 \sum_{\substack{\pi:H\hookrightarrow[r]\\(i,\pi(i))\in\Omega}}
       \left(\prod_{i\in H}a_{i,\pi(i)}\right)
       \operatorname{per}\bigl(A_{[r]\setminus H,\,[r]\setminus\pi(H)}\bigr).
\]
This is simply the partition of the permutations in the
permanent definition according to their restrictions to
$H$; each full permutation occurs in exactly one summand.
There is no sign and no multiplicity factor.  If $H$ is
empty the sum has its one empty restriction, and its
weight is one.  Enumerating at most $r^h$ functions and
retaining the allowed injections constructs the sum
deterministically with polynomial overhead per term.

Every minor has order $r'=r-h\ge r/2$.  Its remaining
row supports have size at most $T$, and deleting columns
cannot increase them.  Since
$\lceil\log_2r'\rceil\ge\lambda-1$, one has
\[
                   T\le(\lceil\log_2r'\rceil)^{a+3}
\]
for sufficiently large $r$.  The preceding row-scope
permanent construction therefore applies to every minor,
with a fixed changed scope exponent.  Zero rows can be
detected from the prescribed support and give zero terms.
The finitely many exceptional matrix orders, including
orders zero and one, are computed directly.

The logarithm of the number of restrictions is at most
\[
                   h\log_2r=O_C(r/\log r),
\]
which is absorbed by
$O_{C,a}(r(\log\log r)^2/\sqrt{\log r})$.
The minor construction has that same bound because
$r'=\Theta(r)$.  Multiplication by the $h$ selected entry
variables and summation of all terms add only the stated
overhead.  Each minor circuit has formal degree at most
$2r'$, so each term has formal degree at most
$2r'+h\le2r$, and summation takes the maximum.  All scalar
constants retain their polynomial bit lengths.  This
proves the total-support theorem as an integral formal
identity, and therefore over every commutative coefficient
ring by base change.

\paragraph{Total-support geometric and representation theorem.}
For a support $\Omega\subseteq[r]\times[r]$ satisfying
$|\Omega|\le Cr(\lceil\log_2r\rceil)^a$, put
$B=\mathbb Z[a_{ij}:(i,j)\in\Omega]$ and
$\mathscr A=B[x_1,\ldots,x_r]/(x_j^2-x_j)$.  The regular
function
\[
 \operatorname{Tr}_{\mathscr A/B}\left(
     (-1)^r\prod_j(1-2x_j)
                 \prod_i\sum_{j:(i,j)\in\Omega}a_{ij}x_j\right)
                              =\operatorname{per}(A)
\]
has an integral constant-free circuit of size
$\operatorname{poly}(r)2^{O_{C,a}(r(\log\log r)^2/\sqrt{\log r})}$
and formal degree at most $2r$, valid after every base
change.  Ryser's identity identifies this signed trace,
while the heavy-row decomposition constructs it using
smaller trace functions with polylogarithmic row scopes.

For the split torus with characters indexed by $\Omega$,
let $\mathcal V_\Omega$ have one basis vector of character
$\prod_i a_{i,\sigma(i)}$ per allowed perfect matching.
Its effective character is the same permanent polynomial.
The restriction partition above gives the actual direct-sum
decomposition
\[
 \mathcal V_\Omega\simeq
  \bigoplus_{\pi}
    \chi_\pi\otimes\mathcal V_{\Omega_\pi},
 \qquad \chi_\pi=\prod_{i\in H}a_{i,\pi(i)},
\]
where $\Omega_\pi$ is the remaining minor support and its
characters are included in the original torus.  Applying
the integral minor circuits to these summands gives the
stated character circuit; positivity of the decomposition
does not remove the virtual differences inside the minor
circuits.  The precise classical trace and character-ring
interfaces remain EGA IV4, 18.2.1--18.2.2, printed p.~111,
and the SGA3/SGA6 locators already stated in NN-S019.
